\chapter{Introdu\c{c}\~ao}
\label{cap:introducao}

Na primeira p\'agina da parte textual, apresente o seu trabalho. A
introdu\c{c}\~ao deve ocupar uma ou, no m\'aximo, duas p\'aginas. Nela,
dê um panorama da tem\'atica, exponha a quest\~ao-problema que conduzir\'a
a pesquisa proposta, a hip\'otese e a justificativa.

Os t\'itulos de cada se\c{c}\~ao (1 INTRODU\c{C}\~AO, 2 REFERENCIAL
TE\'ORICO etc.) devem ser grafados em letras mai\'usculas e negritadas,
sem ponto ap\'os a numera\c{c}\~ao (a classe \texttt{idpthesis} j\'a faz
isso automaticamente para \verb|\chapter|).

\section{Objetivo geral}
Descreva aqui o objetivo geral do trabalho.

\section{Objetivos espec\'ificos}
\begin{itemize}
  \item Objetivo espec\'ifico 1;
  \item Objetivo espec\'ifico 2.
\end{itemize}

\section{Problema}
Apresente a quest\~ao-problema que conduz a pesquisa.

\section{Hip\'otese}
Apresente a hip\'otese de pesquisa.

\section{Justificativa}
Justifique a relev\^ancia do trabalho.

Caso haja necessidade de notas explicativas, use notas de rodap\'e\footnote{Texto
de exemplo para nota de rodap\'e. A nota de rodap\'e apresenta informa\c{c}\~oes
ou coment\'arios adicionais sobre o conte\'udo.}, que a classe já formata em
fonte tamanho 10, conforme a norma.
